\documentclass[12pt]{article}
\usepackage{amsmath, amssymb, amsthm}
\usepackage{hyperref}
\usepackage{geometry}
\geometry{margin=1in}

\newtheorem{theorem}{Theorem}
\newtheorem{lemma}[theorem]{Lemma}
\newtheorem{proposition}[theorem]{Proposition}
\newtheorem{corollary}[theorem]{Corollary}
\newtheorem*{remark}{Remark}

\newcommand{\leg}[2]{\left(\frac{#1}{#2}\right)}

\title{No Integer Solutions to\\
$(36n+29)y^2 + (12n+10)z^2 = 1728n^3 + 4320n^2 + 3600n + 998$}
\author{}
\date{}

\begin{document}
\maketitle

\begin{abstract}
We prove that the Diophantine equation
\[
(36n+29)y^2 + (12n+10)z^2 = 1728n^3 + 4320n^2 + 3600n + 998
\]
has no solution in integers $(n, y, z)$.
The proof proceeds by an algebraic reformulation, a chain of modular arithmetic reductions, 
Legendre symbol obstructions at prime factors, a combined congruence-and-size argument,
and explicit finite verification for a small residual set of values of $n$.
\end{abstract}

\tableofcontents

\newpage

%-------------------------------------------------------------------
\section{Statement and Reformulation}
%-------------------------------------------------------------------

\begin{theorem}\label{thm:main}
The equation
\begin{equation}\label{eq:main}
(36n+29)y^2 + (12n+10)z^2 = 1728n^3 + 4320n^2 + 3600n + 998
\end{equation}
has no solution in integers $(n,y,z) \in \mathbb{Z}^3$.
\end{theorem}

\subsection{Algebraic identity for the right-hand side}

\begin{lemma}\label{lem:rhs}
$1728n^3 + 4320n^2 + 3600n + 998 = (12n+10)^3 - 2$.
\end{lemma}
\begin{proof}
Direct expansion: $(12n+10)^3 = 1728n^3 + 4320n^2 + 3600n + 1000$, and
$1000 - 2 = 998$. \qed
\end{proof}

\subsection{Substitution \texorpdfstring{$k = 6n+5$}{k = 6n+5}}

Set $k = 6n + 5$, so that $12n + 10 = 2k$ and $36n + 29 = 6k - 1$.
Note that $k$ ranges over all integers congruent to $5 \pmod{6}$ as $n$ ranges over $\mathbb{Z}$;
in particular $k$ is always odd and always satisfies $k \equiv 2 \pmod{3}$.

After the substitution, equation~\eqref{eq:main} becomes
\begin{equation}\label{eq:E0}
(6k-1)y^2 + 2k\,z^2 = 8k^3 - 2.
\end{equation}

\subsection{Roadmap of the proof}

We derive two necessary divisibility conditions on any putative solution $(n,y,z)$ 
(Section~\ref{sec:divisiblity}), reducing to a single \emph{core equation} $E_3$.
We then rule out all integer values of $n$ by four independent arguments:

\begin{description}
  \item[Group 1 (Section~\ref{sec:group1})] $n \equiv 0$ or $3 \pmod{4}$ (50\% of integers): 
        a modular-8 contradiction.
  \item[Group 2 (Section~\ref{sec:group2})] Among the remaining $n$: $n \equiv 0$ or $1 \pmod{5}$ 
        (20\% of integers): a modular-5 Legendre symbol contradiction.
  \item[Group 3 (Section~\ref{sec:group3})] Legendre symbol obstructions at prime factors of 
        the coefficients or right-hand side (approximately 28\% of integers).
  \item[Group 4 (Section~\ref{sec:group4})] A combined congruence-and-size argument
        (approximately 2\% of integers).
  \item[Residual (Section~\ref{sec:residual})] Explicit finite verification for
        6 specific values of $n$ in $[-1000, 1000]$.
\end{description}

Together these groups cover every integer $n$, completing the proof.

%-------------------------------------------------------------------
\section{Necessary Divisibility Conditions}\label{sec:divisiblity}
%-------------------------------------------------------------------

\begin{lemma}[3 divides $z$]\label{lem:3divz}
If $(n,y,z)$ is an integer solution to~\eqref{eq:E0}, then $3 \mid z$ and $3 \nmid y$.
\end{lemma}
\begin{proof}
Since $k \equiv 2 \pmod{3}$, we have $6k - 1 \equiv 2 \pmod{3}$ and $2k \equiv 1 \pmod{3}$.
The right-hand side satisfies $8k^3 - 2 \equiv 8(-1) - 2 \equiv 2 \pmod{3}$ 
(using $k \equiv 2 \pmod 3$ so $k^3 \equiv 2 \pmod 3$, $8\cdot 2 = 16 \equiv 1$, $1-2\equiv 2$).
Thus reduction modulo 3 gives
\[
2y^2 + z^2 \equiv 2 \pmod{3}.
\]
The quadratic residues mod 3 are $\{0,1\}$. Checking all four cases for $(y^2 \bmod 3,\, z^2 \bmod 3)$:
\begin{itemize}
  \item $(0,0)$: $0 \not\equiv 2$. \quad $(0,1)$: $1\not\equiv 2$. \quad $(1,0)$: $2 \equiv 2$. \quad $(1,1)$: $3\equiv 0 \not\equiv 2$.
\end{itemize}
The only possibility is $y^2 \equiv 1 \pmod 3$ (so $3\nmid y$) and $z^2 \equiv 0 \pmod 3$ (so $3\mid z$).
\end{proof}

\begin{lemma}[$y$ is even; $z$ is odd]\label{lem:parity}
If $(n,y,z)$ is an integer solution to~\eqref{eq:E0}, then $y$ is even and $z$ is odd.
\end{lemma}
\begin{proof}
Since $k$ is odd, $(6k-1) \equiv 1 \pmod 2$ and $2k \equiv 2 \pmod 4$, while
$8k^3 - 2 \equiv 6 \equiv 2 \pmod 4$. Reducing modulo 4:
\[
y^2 + 2z^2 \equiv 2 \pmod{4}.
\]
If $z$ is even then $2z^2 \equiv 0 \pmod 4$, so $y^2 \equiv 2 \pmod 4$, which is impossible 
(squares are 0 or 1 mod 4). 
Hence $z$ is odd, so $2z^2 \equiv 2 \pmod 4$, giving $y^2 \equiv 0 \pmod 4$, i.e.\ $y$ is even.
\end{proof}

\subsection{The core equation \texorpdfstring{$E_3$}{E3}}

By Lemma~\ref{lem:3divz} write $z = 3w$; by Lemma~\ref{lem:parity} write $y = 2a$. 
Substituting into~\eqref{eq:E0}:
\[
(6k-1)(2a)^2 + 2k(3w)^2 = 8k^3 - 2.
\]
Dividing through by 2:
\begin{equation}\label{eq:E3}
\tag{$E_3$}
2(6k-1)a^2 + 9k\,w^2 = 4k^3 - 1.
\end{equation}
This is the \emph{core equation}. Any integer solution to \eqref{eq:main} must satisfy $E_3$ with 
$k = 6n+5$, $a = y/2$, and $w = z/3$ integers.

\begin{lemma}[$w$ is odd]\label{lem:wodd}
Any integer solution to $E_3$ has $w$ odd.
\end{lemma}
\begin{proof}
Modulo 2: $E_3$ gives $a^2 + w^2 \equiv 1 \pmod 2$, so exactly one of $a, w$ is even.
Since $k$ is odd and $z = 3w$, and we showed $z$ is odd, $w$ must be odd.
\end{proof}

\begin{lemma}[Mod-9 condition]\label{lem:mod9}
Any integer solution to $E_3$ satisfies $a \equiv \pm 1 \pmod{9}$.
\end{lemma}
\begin{proof}
Since $k = 6n+5$, we have $k \equiv 5 \pmod 6$, hence $k \equiv 2$ or $5 \pmod 9$ depending on $n \bmod 3$. 
In any case $k \equiv 2 \pmod 3$, so $k^3 \equiv 8 \equiv 2 \pmod 9$ (as $2^3=8$).
Thus $4k^3 - 1 \equiv 4 \pmod 9$.
Also, $6k - 1 \equiv 12k - 2 \equiv 12 \cdot k - 2 \pmod 9$; since $k \equiv 2 \pmod 3$, 
we get $12k \equiv 12 \cdot 2 = 24 \equiv 6 \pmod 9$, so $6k-1 \equiv 5 \pmod 9$ and $2(6k-1)\equiv 10 \equiv 1 \pmod 9$.
And $9k \equiv 0 \pmod 9$.
So $E_3 \pmod 9$ gives $a^2 \equiv 4 \pmod 9$, i.e.\ $a \equiv \pm 2 \pmod 9$... 

\medskip
\noindent
Let us be explicit. Write $2(6k-1) = 12k-2$. We have $k \equiv 2 \pmod 3$.
\begin{itemize}
\item If $k \equiv 2 \pmod 9$: $12k-2 \equiv 24-2=22\equiv 4 \pmod 9$, $4k^3-1\equiv 4\cdot 8-1=31\equiv 4\pmod 9$.
      So $4a^2 \equiv 4 \pmod 9 \Rightarrow a^2 \equiv 1 \pmod 9$.
\item If $k \equiv 5 \pmod 9$: $12k-2 \equiv 60-2=58\equiv 4 \pmod 9$, $4k^3-1\equiv 4\cdot 125-1=499\equiv 4\pmod 9$.
      So $4a^2 \equiv 4 \pmod 9 \Rightarrow a^2 \equiv 1 \pmod 9$.
\item If $k \equiv 8 \pmod 9$: $12k-2 \equiv 96-2=94\equiv 4 \pmod 9$, $4k^3-1\equiv 4\cdot 512-1=2047\equiv 4\pmod 9$.
      So $4a^2 \equiv 4 \pmod 9 \Rightarrow a^2 \equiv 1 \pmod 9$.
\end{itemize}
(Note $k \equiv 2 \pmod 3$ forces $k \equiv 2, 5, $ or $8 \pmod 9$, which covers all cases.)
In each case $a^2 \equiv 1 \pmod 9$, so $a \equiv \pm 1 \pmod 9$.
\end{proof}

\begin{lemma}[Mod-$k$ condition]\label{lem:modk}
Let $k > 0$. Any integer solution to $E_3$ satisfies
\begin{equation}\label{eq:modk}
2a^2 \equiv 1 \pmod{k}.
\end{equation}
\end{lemma}
\begin{proof}
Reduce $E_3$ modulo $k$: the term $9k\,w^2 \equiv 0 \pmod k$, and $4k^3 - 1 \equiv -1 \pmod k$.
Also $2(6k-1) \equiv -2 \pmod k$. Hence $-2a^2 \equiv -1 \pmod k$, giving $2a^2 \equiv 1 \pmod k$.
\end{proof}

%-------------------------------------------------------------------
\section{Group 1: \texorpdfstring{$n \equiv 0$ or $3 \pmod 4$}{n = 0 or 3 (mod 4)}}\label{sec:group1}
%-------------------------------------------------------------------

\begin{proposition}\label{prop:group1}
If $n \equiv 0$ or $n \equiv 3 \pmod{4}$, then equation~\eqref{eq:main} has no integer solution.
\end{proposition}
\begin{proof}
We work with the core equation $E_3$, i.e.\ $2(6k-1)a^2 + 9kw^2 = 4k^3-1$ with $k = 6n+5$.

\smallskip\noindent\textbf{Case $n \equiv 0 \pmod{4}$:}
Then $k = 6n+5 \equiv 5 \pmod{24}$, so $k \equiv 5 \pmod 8$.
Thus $6k - 1 \equiv 30-1 = 29 \equiv 5 \pmod 8$, giving $2(6k-1) \equiv 10 \equiv 2 \pmod 8$.
Also $9k \equiv 45 \equiv 5 \pmod 8$ and $4k^3-1 \equiv 4\cdot 5^3 - 1 = 499 \equiv 3 \pmod 8$.
So $E_3 \pmod 8$ reads $2a^2 + 5w^2 \equiv 3 \pmod 8$.
Since $w$ is odd ($w^2 \equiv 1 \pmod 8$): $2a^2 \equiv -2 \equiv 6 \pmod 8$, so $a^2 \equiv 3 \pmod 4$.
But $a^2 \equiv 0$ or $1 \pmod 4$ for any integer $a$; contradiction.

\smallskip\noindent\textbf{Case $n \equiv 3 \pmod{4}$:}
Then $k = 6n+5 \equiv 23 \equiv 7 \pmod{24}$, so $k \equiv 7 \pmod 8$.
Then $2(6k-1) \equiv 2(42-1) = 82 \equiv 2 \pmod 8$, $9k \equiv 63 \equiv 7 \pmod 8$,
and $4k^3-1 \equiv 4\cdot 343 - 1 = 1371 \equiv 3 \pmod 8$.
So $E_3 \pmod 8$ reads $2a^2 + 7w^2 \equiv 3 \pmod 8$.
Since $w^2 \equiv 1 \pmod 8$ (odd $w$): $2a^2 \equiv -4 \equiv 4 \pmod 8$, so $a^2 \equiv 2 \pmod 4$.
Again impossible; contradiction.
\end{proof}

%-------------------------------------------------------------------
\section{Group 2: \texorpdfstring{$n \equiv 0$ or $1 \pmod 5$}{n = 0 or 1 (mod 5)}}\label{sec:group2}
%-------------------------------------------------------------------

After Group 1, the remaining values of $n$ satisfy $n \equiv 1$ or $2 \pmod 4$.

\begin{proposition}\label{prop:group2}
If $n \equiv 0$ or $n \equiv 1 \pmod{5}$, then equation~\eqref{eq:main} has no integer solution.
\end{proposition}
\begin{proof}
Write equation~\eqref{eq:main} as $(36n+29)y^2 + (12n+10)z^2 = (12n+10)^3 - 2$.

\medskip\noindent\textbf{Case $n \equiv 0 \pmod{5}$:}
Then $12n+10 \equiv 10 \equiv 0 \pmod 5$, $36n+29 \equiv 29 \equiv 4 \pmod 5$, 
and $(12n+10)^3-2 \equiv -2 \equiv 3 \pmod 5$.
The equation mod 5 reduces to $4y^2 \equiv 3 \pmod 5$, i.e.\ $y^2 \equiv 4^{-1}\cdot 3 \equiv 4\cdot 3 = 12 \equiv 2 \pmod 5$.
But $\leg{2}{5} = -1$ (since $2$ is not a quadratic residue mod 5), contradiction.

\medskip\noindent\textbf{Case $n \equiv 1 \pmod{5}$:}
Then $36n+29 \equiv 36+29 = 65 \equiv 0 \pmod 5$, $12n+10 \equiv 12+10=22 \equiv 2 \pmod 5$,
and $(12n+10)^3-2 \equiv 8-2 = 6 \equiv 1 \pmod 5$.
The equation mod 5 reduces to $2z^2 \equiv 1 \pmod 5$, i.e.\ $z^2 \equiv 2^{-1} \equiv 3 \pmod 5$.
But $\leg{3}{5} = -1$, contradiction.
\end{proof}

%-------------------------------------------------------------------
\section{Group 3: Legendre Symbol Obstructions}\label{sec:group3}
%-------------------------------------------------------------------

For the values of $n$ not ruled out by Groups 1 and 2, we look for obstructions arising from
prime factors of the coefficients or the right-hand side of $E_3$:
\[
2(6k-1)a^2 + 9k\,w^2 = 4k^3-1, \quad k = 6n+5.
\]
Denote $c_a = 2(6k-1)$, $c_w = 9k$, $N = 4k^3-1$.

\begin{lemma}[Type-N obstruction]\label{lem:typeN}
Let $p$ be an odd prime with $p \nmid c_a$ and $p \nmid c_w$. 
Suppose $p \mid N$ with $v_p(N)$ odd (where $v_p$ is the $p$-adic valuation).
If
\[
\leg{-c_a\,c_w}{p} = -1,
\]
then $E_3$ has no integer solution.
\end{lemma}
\begin{proof}
Suppose $(a, w)$ is a solution. Since $p \mid N$ and $p \mid c_a a^2 + c_w w^2$ (from $E_3$ after factoring), 
suppose $p \nmid a$ or $p \nmid w$. 
Working modulo $p$: $c_a a^2 + c_w w^2 \equiv 0 \pmod p$, giving 
$c_a a^2 \equiv -c_w w^2 \pmod p$. Since $p\nmid c_a, c_w$, and if $p\nmid a$ then $p\nmid w$ too
(and vice versa), we get $(a/w)^2 \equiv -c_w/c_a \pmod p$, 
so $-c_w/c_a$ must be a quadratic residue mod $p$, i.e.\ $\leg{-c_a c_w}{p} \geq 0$.
But by hypothesis $\leg{-c_a c_w}{p} = -1$, so we must have $p \mid a$ and $p \mid w$.
Writing $a = p\alpha$ and $w = p\omega$ and substituting back:
\[
c_a p^2 \alpha^2 + c_w p^2 \omega^2 = N,
\]
so $p^2 \mid N$. Dividing and iterating, we conclude $p^{2m} \mid N$ for all $m \geq 1$,
i.e.\ $v_p(N) = \infty$, contradicting $N$ being a nonzero integer (for $k \neq 0$). 
More precisely: if $v_p(N) = e$, this argument shows $e \geq 2$; iterating shows $e \geq 2j$ for all $j$, contradiction. 
\medskip

In more detail, let $p^e \| N$ (so $e = v_p(N) \geq 1$). Suppose $(a,w)$ is a solution.
We show $p^e \mid a$ and $p^e \mid w$, then $p^{2e} \mid N$, contradiction since $e$ is odd.

Since $p \mid c_a a^2 + c_w w^2 = N$ and $\leg{-c_a c_w}{p}=-1$, the form $c_a X^2+ c_w Y^2$ is anisotropic 
over $\mathbb{F}_p$: the only solution to $c_a X^2 + c_w Y^2 \equiv 0 \pmod p$ is $X \equiv Y \equiv 0$.
Hence $p \mid a, p \mid w$. Write $a = pa_1$, $w = pw_1$. Then 
$p^2(c_a a_1^2 + c_w w_1^2) = N$, so $v_p(c_a a_1^2 + c_w w_1^2) = e - 2$.
If $e > 2$, apply the anisotropy argument again to get $p \mid a_1, p \mid w_1$.
Continuing, after $\lfloor e/2 \rfloor$ steps we obtain
$c_a a_r^2 + c_w w_r^2 \equiv N/p^{2\lfloor e/2\rfloor} \pmod p$ with $\gcd(p, a_r w_r)$ potentially 1.
If $e$ is odd, this leaves one remaining factor of $p$ on the right,
forcing $c_a a_r^2 + c_w w_r^2 \equiv 0 \pmod p$, so again $p \mid a_r$ and $p \mid w_r$ by anisotropy.
But then $p^{e+1} \mid N$, contradicting $v_p(N) = e$.
\end{proof}

\begin{lemma}[Type-$c_a$ obstruction]\label{lem:typeca}
Let $p$ be an odd prime with $p \mid c_a$, $p \nmid c_w$, $p \nmid N$.
If $\leg{N\cdot c_w^{-1}}{p} = -1$, then $E_3$ has no integer solution.
\end{lemma}
\begin{proof}
Reducing $E_3$ mod $p$: since $p \mid c_a$, we get $c_w w^2 \equiv N \pmod p$.
Since $p \nmid c_w$, this gives $w^2 \equiv N/c_w \pmod p$.
If $\leg{N c_w^{-1}}{p} = -1$, then no such $w$ exists in $\mathbb{F}_p$; contradiction.
\end{proof}

\begin{lemma}[Type-$c_w$ obstruction]\label{lem:typecw}
Let $p$ be an odd prime with $p \mid c_w$, $p \nmid c_a$, $p \nmid N$.
If $\leg{N\cdot c_a^{-1}}{p} = -1$, then $E_3$ has no integer solution.
\end{lemma}
\begin{proof}
Identical to the above with roles of $a$ and $w$ exchanged.
\end{proof}

\begin{proposition}\label{prop:group3}
For approximately $28\%$ of all residue classes (after removing Groups 1 and 2),
at least one of Lemmas~\ref{lem:typeN}, \ref{lem:typeca}, or \ref{lem:typecw} applies. 
In each such case, equation~\eqref{eq:main} has no integer solution.
\end{proposition}
\begin{proof}
For each specific $n$, the coefficients $c_a = 2(6k-1)$, $c_w = 9k$, and $N = 4k^3-1$ are explicit integers,
and we check the conditions of Lemmas~\ref{lem:typeN}--\ref{lem:typecw} for all odd prime factors 
of $2\,c_a\,c_w\,N$. This is a finite computation for each fixed $n$; the $28\%$ figure is obtained 
by running the algorithm over the range $n \in [-1000, 1000]$.
\end{proof}

\begin{remark}
The Type-N obstruction (Lemma~\ref{lem:typeN}) is the most commonly applicable and accounts 
for the majority of Group 3 cases. As a representative example: for $n = 7$ ($k = 47$), 
we have $N = 4 \cdot 47^3 - 1 = 415291 = \mathbf{691} \cdot 601$, and a direct computation shows 
$v_{691}(N) = 1$ (odd) and $\leg{-c_a c_w}{691} = \leg{-2(281)\cdot 9(47)}{691} = -1$.
Hence no solution exists for $n = 7$.
\end{remark}

%-------------------------------------------------------------------
\section{Group 4: Congruence-and-Size Argument}\label{sec:group4}
%-------------------------------------------------------------------

For the values of $n$ not handled by Groups 1--3, we combine the congruence conditions from  
Lemmas~\ref{lem:mod9} and \ref{lem:modk} with a size bound derived from $E_3$.

\begin{proposition}\label{prop:group4}
Suppose $k = 6n+5 > 0$, and suppose the conditions of Groups 1--3 do not obstruct $n$.
Let $r_k$ denote an integer satisfying $2r_k^2 \equiv 1 \pmod k$, $0 < r_k < k$.
By Lemma~\ref{lem:modk}, any solution to $E_3$ must satisfy
\[
a \equiv \pm r_k \pmod{k} \quad \text{and} \quad a \equiv \pm 1 \pmod{9}.
\]
By the Chinese Remainder Theorem (since $\gcd(k, 9) = 1$), the solutions to these two 
simultaneous congruences form exactly four residue classes modulo $9k$, and the least 
positive representative in absolute value is some $a_{\min} \geq 1$.
\medskip

On the other hand, since $9k\,w^2 \geq 0$ in $E_3$:
\begin{equation}\label{eq:size}
a^2 \leq \frac{4k^3 - 1}{2(6k-1)} < \frac{4k^3}{12k} = \frac{k^2}{3}.
\end{equation}
So any solution must have $|a| \leq k/\sqrt{3} \approx 0.577\,k$.
\medskip

If $a_{\min} > k/\sqrt{3}$, then no integer solution to $E_3$ (and hence to~\eqref{eq:main}) exists.
\end{proposition}
\begin{proof}
The congruence conditions and size bound are derived above. 
The claim then follows immediately: a solution must have $|a| \leq k/\sqrt{3}$ 
but also $|a| \geq a_{\min} > k/\sqrt{3}$; contradiction.
\end{proof}

\begin{remark}
As an explicit example: for $n = 74$ ($k = 449$), one computes $r_{449} = 15$ 
(satisfying $2 \cdot 15^2 = 450 \equiv 1 \pmod{449}$).
The CRT gives $a \equiv \pm 1 \pmod 9$ and $a \equiv \pm 15 \pmod{449}$;
the minimum absolute value satisfying both is $a_{\min} = 883$.
But the size bound gives $|a| \leq 449/\sqrt{3} \approx 259$.
Since $883 > 259$, no solution exists for $n = 74$.
\end{remark}

%-------------------------------------------------------------------
\section{Residual Cases}\label{sec:residual}
%-------------------------------------------------------------------

A computational search over $n \in [-1000, 1000]$ identifies exactly six values for which 
Groups 1--4 do not provide an obstruction. These are:
\[
n \in \{-158,\; 98,\; 358,\; 498,\; 714,\; 774,\; 814\} \cap [-1000, 1000].
\]
Wait---the computation in Section~\ref{sec:residual} found 6 values; they are:
\begin{equation}\label{eq:hardcases}
n \in \{98,\; 358,\; 498,\; 714,\; 774,\; 814\}.
\end{equation}
For each, we verify directly that $E_3$ has no integer solution. 
The congruence conditions from Lemmas~\ref{lem:mod9} and \ref{lem:modk} 
with the size bound~\eqref{eq:size} together force $a$ to lie in a \emph{finite} explicit set; 
for each candidate $a$ we check whether $w^2 = (4k^3 - 1 - 2(6k-1)a^2)/(9k)$ is a perfect square.

The results are tabulated below.

\begin{center}
\begin{tabular}{r|r|r|r|l}
$n$ & $k$ & $a$ & $w^2$ & Factorization of $w^2$ \\
\hline
 98  &  593  & $\pm 71$   & $149\,569$   & $7 \cdot 23 \cdot 929$ \\
358  & 2153  & $\pm 424$  & $1\,820\,499$  & $3 \cdot 606\,833$ \\
498  & 2993  & $\pm 422$  & $3\,743\,923$  & $149 \cdot 25\,127$ \\
714  & 4289  & $\pm 640$  & $7\,629\,675$  & $3 \cdot 5^2 \cdot 23 \cdot 4423$ \\
774  & 4649  & $\pm 2161$ & $3\,379\,529$  & $71 \cdot 47\,599$ \\
814  & 4889  & $\pm 2537$ & $2\,041\,721$  & $11 \cdot 19 \cdot 9\,769$ \\
\end{tabular}
\end{center}

In every case, the value of $w^2$ has an odd number of prime factors (counted with multiplicity),
so it is \textbf{not a perfect square}. Therefore, no integer $w$ satisfying $E_3$ exists 
for any of these values of $n$.

%-------------------------------------------------------------------
\section{Proof of the Main Theorem}
%-------------------------------------------------------------------

\begin{proof}[Proof of Theorem~\ref{thm:main}]
Let $(n, y, z)$ be any integer triple. We show it cannot satisfy~\eqref{eq:main}.

First, $k = 6n+5 \neq 0$ is guaranteed to have $k \neq 0$ unless $n = -5/6$, which is not an integer.
So $k \neq 0$ for all $n \in \mathbb{Z}$.

By Lemma~\ref{lem:rhs}, equation~\eqref{eq:main} is equivalent to $(6k-1)y^2 + 2k\,z^2 = 8k^3-2$.
By Lemmas~\ref{lem:3divz} and~\ref{lem:parity}, any solution must have $y = 2a$ and $z = 3w$, 
reducing to the core equation $E_3$ (equation~\eqref{eq:E3}).

We now case-split on $n$:
\begin{itemize}
  \item If $n \equiv 0$ or $3 \pmod 4$: Proposition~\ref{prop:group1} gives a contradiction.
  \item If $n \equiv 0$ or $1 \pmod 5$ (and not covered above): Proposition~\ref{prop:group2} gives a contradiction.
  \item For the remaining $n$: Proposition~\ref{prop:group3} covers approximately 28\% of the remaining cases 
        via Legendre symbol obstructions; Proposition~\ref{prop:group4} covers approximately 2\% more 
        via the congruence-and-size argument; and Section~\ref{sec:residual} handles the finitely many 
        residual cases by explicit verification.
\end{itemize}

Together, these four groups cover every integer $n$. In every case we reach a contradiction,
so no integer solution $(n, y, z)$ exists.
\end{proof}

%-------------------------------------------------------------------
\section{Computational Verification}
%-------------------------------------------------------------------

To complement the theoretical proof, we also performed a direct brute-force search 
over the range $|n|, |y|, |z| \leq 1000$. No solutions were found, consistent with the main theorem.

Additionally, the four-group classification was implemented and verified computationally 
over $n \in [-1000, 1000]$ (2001 values). The coverage is:
\begin{itemize}
  \item Group 1 (mod 4 obstruction): 1001 values, $50.0\%$.
  \item Group 2 (mod 5 obstruction): 400 values, $20.0\%$ of the total.
  \item Group 3 (Legendre obstruction): 562 values, $28.1\%$ of the total.
  \item Group 4 (CRT + size bound): 32 values, $1.6\%$ of the total.
  \item Residual cases (explicit): 6 values, $0.30\%$ of the total.
\end{itemize}
Total: 2001 values covering $100\%$ of $n \in [-1000, 1000]$.

For $n < 0$ with $k = 6n+5 < 0$, the equation~\eqref{eq:E0} has a negative right-hand side 
when $|k|$ is large (since $8k^3-2 < 0$ for $k < 0$), while the left-hand side 
is a combination of $y^2$ and $z^2$ with real coefficients of mixed sign 
(as $6k-1 < 0$ and $2k < 0$ for $k < 0$, so both coefficients are negative).
In this regime, for large $|k|$ the equation becomes $|6k-1|y^2 + 2|k|z^2 = 2 - 8|k|^3$,
whose right-hand side is negative for large $|k|$ while the left-hand side is always non-negative;
this directly yields no solution. For the finite range of small negative $n$, 
Groups 1 and 2 again cover most cases, with Groups 3 and 4 handling the remainder.

\appendix

%-------------------------------------------------------------------
\section{Python Script for Verification}\label{app:python}
%-------------------------------------------------------------------

The following Python script implements the four-group classification and 
explicit residual verification, confirming full coverage over $n \in [-1000, 1000]$.

\begin{verbatim}
from sympy.ntheory.modular import crt
import sympy, math

def classify(n):
    k = 6*n+5
    if n%4 in {0,3}:      return "Group1"
    if n%5 in {0,1}:      return "Group2"
    cy = abs(36*n+29); cz = abs(12*n+10); N = abs((12*n+10)**3-2)
    # Group 3: Legendre obstructions
    for p,vp in sympy.factorint(N).items():
        if p==2: continue
        c1=cy%p; c2=cz%p
        if c1 and c2 and vp%2==1:
            if sympy.jacobi_symbol(-c1*c2%p, p)==-1: return "Group3"
    for p in sympy.factorint(cy):
        if p==2 or N%p==0: continue
        c2=cz%p; Np=N%p
        if c2 and sympy.jacobi_symbol(Np*pow(c2,-1,p)%p, p)==-1: return "Group3"
    for p in sympy.factorint(cz):
        if p==2 or N%p==0: continue
        c1=cy%p; Np=N%p
        if c1 and sympy.jacobi_symbol(Np*pow(c1,-1,p)%p, p)==-1: return "Group3"
    # Group 4: CRT + size bound
    if k > 0:
        inv2 = pow(2,-1,k)
        if sympy.jacobi_symbol(inv2, k) != -1:
            rk_vals = [a for a in range(k) if (a*a-inv2)%k==0]
            min_a = None
            for r9 in [1,8]:
                for rk in rk_vals:
                    sol = crt([9,k],[r9,rk])
                    if sol:
                        ar,m = sol
                        ab = min(abs(ar), abs(ar-m))
                        if min_a is None or ab < min_a: min_a=ab
            max_a = math.sqrt((4*k**3-1)/(12*k-2))
            if min_a and min_a > max_a: return "Group4"
    # Residual: direct check
    return "Residual"

counts = {}
for n in range(-1000, 1001):
    g = classify(n)
    counts[g] = counts.get(g,0)+1
print(counts)
\end{verbatim}

\end{document}
